\vspace*{20pt}

\section{Polímeros o macromoléculas}

\subsection{¿Qué es un polímero?}

Los \textbf{polímeros} son macromoléculas formadas por la unión repetitiva de unidades
más pequeñas denominadas \textbf{monómeros}, enlazadas mediante enlaces covalentes.
Su masa molecular puede alcanzar millones de unidades de masa atómica, por lo que
presentan propiedades muy diferentes a las de las moléculas pequeñas.

La palabra polímero procede del griego \emph{poly} (muchos) y \emph{meros} (partes).

\begin{center}
\[
n\,\mathrm{M}\longrightarrow (\mathrm{M})_n
\]
\end{center}

Donde $M$ representa el monómero y $n$ el grado de polimerización.

\subsection{Estructura}

Las cadenas poliméricas pueden ser muy largas y establecer interacciones entre ellas
(fuerzas de Van der Waals, puentes de hidrógeno o enlaces cruzados). Estas
interacciones son responsables de propiedades como la rigidez, elasticidad o
temperatura de fusión.

\begin{center}
\begin{tabular}{c}
Monómero $\rightarrow$ Cadena polimérica $\rightarrow$ Material\\
CH$_2$=CH$_2$ $\rightarrow$ (-CH$_2$-CH$_2$-)$_n$ $\rightarrow$ Polietileno
\end{tabular}
\end{center}

\subsection{Características generales}

\begin{itemize}
\item Elevada masa molecular.
\item Baja densidad en comparación con los metales.
\item Buenos aislantes eléctricos y térmicos.
\item Gran resistencia química.
\item Facilidad para moldearse.
\end{itemize}

\subsection{Polímeros naturales y sintéticos}

Los polímeros aparecen tanto en la naturaleza como fabricados industrialmente.

\begin{center}
% --- tabla antigua ---
% \begin{tabular}{p{0.42\linewidth}p{0.42\linewidth}}
% \textbf{Naturales} & \textbf{Sintéticos}\\
% \hline
% Celulosa & Polietileno (PE)\\
% Almidón & PVC\\
% Proteínas & Polipropileno (PP)\\
% ADN & PET\\
% Caucho natural & Nailon\\
% \end{tabular}
% --- tabla nueva ---
\begin{tabular}{@{}p{0.42\linewidth}p{0.42\linewidth}@{}}
    \toprule
    \textbf{Naturales} & \textbf{Sintéticos} \\
    \midrule
    Celulosa & Polietileno (PE) \\[3pt]
    Almidón & PVC \\[3pt]
    Proteínas & Polipropileno (PP) \\[3pt]
    ADN & PET \\[3pt]
    Caucho natural & Nailon \\
    \bottomrule
\end{tabular}
\end{center}

\begin{tcolorbox}[title=Recuerda,colback=blue!5,colframe=blue]
Un polímero está formado por un gran número de monómeros unidos mediante enlaces
covalentes. Cuanto mayor es el número de monómeros, mayor suele ser la masa molecular
y mejores son sus propiedades mecánicas.
\end{tcolorbox}

\subsection*{Ejemplo}

El eteno es un monómero. Al unirse miles de moléculas de eteno mediante una reacción
de polimerización se obtiene polietileno:

\[
n\,CH_2=CH_2 \longrightarrow (-CH_2-CH_2-)_n
\]

Este plástico se emplea en bolsas, envases y recipientes.

\begin{figure}[H]
\centering

\begin{tikzpicture}[
>=Latex,
font=\small,
node distance=8mm,
monomer/.style={
draw,
rounded corners,
fill=cyan!20,
minimum width=2.2cm,
minimum height=0.9cm,
align=center
},
polymer/.style={
draw,
rounded corners,
fill=orange!20,
minimum width=10cm,
minimum height=1cm,
align=center
}
]

%%%%%%%%%%%%%%%%%%%%%%%%%%%%%%%%%%%%%%%%%%%%%%%%%%%%%%%%%%%%
% Monómeros
%%%%%%%%%%%%%%%%%%%%%%%%%%%%%%%%%%%%%%%%%%%%%%%%%%%%%%%%%%%%

\node[monomer] (m1) {Monómero};

\node[right=of m1] (plus1) {$+$};

\node[monomer,right=of plus1] (m2) {Monómero};

\node[right=of m2] (plus2) {$+$};

\node[monomer,right=of plus2] (m3) {Monómero};

\node[right=of m3] (dots) {$\cdots$};

%%%%%%%%%%%%%%%%%%%%%%%%%%%%%%%%%%%%%%%%%%%%%%%%%%%%%%%%%%%%
% Flecha
%%%%%%%%%%%%%%%%%%%%%%%%%%%%%%%%%%%%%%%%%%%%%%%%%%%%%%%%%%%%

\draw[->,very thick]
($(m2.south)+(0,-0.2)$)
--
node[right]{\textbf{Polimerización}}
++(0,-1.6);

%%%%%%%%%%%%%%%%%%%%%%%%%%%%%%%%%%%%%%%%%%%%%%%%%%%%%%%%%%%%
% Polímero
%%%%%%%%%%%%%%%%%%%%%%%%%%%%%%%%%%%%%%%%%%%%%%%%%%%%%%%%%%%%

\node[
polymer,
below=2.3cm of m2
] (poly)
{
Monómero -- Monómero -- Monómero -- Monómero -- Monómero -- Monómero
};

%%%%%%%%%%%%%%%%%%%%%%%%%%%%%%%%%%%%%%%%%%%%%%%%%%%%%%%%%%%%
% Texto inferior
%%%%%%%%%%%%%%%%%%%%%%%%%%%%%%%%%%%%%%%%%%%%%%%%%%%%%%%%%%%%

\node[
below=5mm of poly
]
{\textbf{Macromolécula (Polímero)}};

%%%%%%%%%%%%%%%%%%%%%%%%%%%%%%%%%%%%%%%%%%%%%%%%%%%%%%%%%%%%
% Flechas laterales
%%%%%%%%%%%%%%%%%%%%%%%%%%%%%%%%%%%%%%%%%%%%%%%%%%%%%%%%%%%%

\draw[->]
(m1.south)
to[out=-90,in=160]
(poly.west);

\draw[->]
(m3.south)
to[out=-90,in=20]
(poly.east);

%%%%%%%%%%%%%%%%%%%%%%%%%%%%%%%%%%%%%%%%%%%%%%%%%%%%%%%%%%%%

\node[
left=2.2cm of poly,
align=center
]
{Muchos\\monómeros};

\node[
right=2.2cm of poly,
align=center
]
{Cadena\\polimérica};

\end{tikzpicture}

\caption{Formación de un polímero por unión repetitiva de monómeros.}

\end{figure}

%%%%%%%%%%%%%%%%%%%%%%%%%%%%%%%%%%%%%%%%%%%%%%%%%%%%%%%%%%%%
%%%%%%%%%%%%%%%%%%%%% SECCIÓN 2 %%%%%%%%%%%%%%%%%%%%%%%%%%%%%
%%%%%%%%%%%%%%%%%%%%%%%%%%%%%%%%%%%%%%%%%%%%%%%%%%%%%%%%%%%%

\section{Clasificación de los polímeros}

Los polímeros pueden clasificarse atendiendo a diferentes criterios. Esta clasificación permite relacionar su estructura con sus propiedades físicas, químicas y sus aplicaciones industriales.

\begin{center}
\begin{tikzpicture}[
    node distance=10mm,
    every node/.style={
        draw,
        rounded corners,
        align=center,
        font=\small,
        minimum width=3cm
    },
    >={Latex}
]

% Nodo principal
\node (P) {Polímeros};

% Primer nivel
\node[below left=of P, xshift=-20mm] (O) {Según\\su origen};
\node[below=of P, xshift=-8mm] (C) {Según\\su composición};
\node[below right=of P, xshift=20mm] (E) {Según\\su estructura};

% Segundo nivel
\node[below=18mm of C, xshift=18mm] (T)
{Según su\\comportamiento\\frente al calor};

% Flechas
\draw[-Latex] (P) -- (O);
\draw[-Latex] (P) -- (C);
\draw[-Latex] (P) -- (E);
\draw[-Latex] (P) -- (T);

\end{tikzpicture}
\end{center}

\vspace{0.3cm}

%%%%%%%%%%%%%%%%%%%%%%%%%%%%%%%%%%%%%%%%%%%%%%%%%%%%%%%%%%%%%

\subsection{Clasificación según su origen}

Dependiendo de su procedencia, los polímeros pueden ser naturales o sintéticos.

\begin{enumerate}[label=\alph*)]

\item \textbf{Polímeros naturales}

Son aquellos que se encuentran en la naturaleza y son sintetizados por los seres vivos.

Algunos ejemplos son:

\begin{itemize}
\item Celulosa.
\item Almidón.
\item Proteínas.
\item ADN y ARN.
\item Caucho natural.
\end{itemize}

Estos materiales suelen ser biodegradables y desempeñan funciones estructurales o biológicas.

\item \textbf{Polímeros sintéticos}

Se obtienen mediante procesos industriales a partir de derivados del petróleo o, cada vez más, de materias primas renovables.

Los más importantes son:

\begin{itemize}
\item Polietileno (PE).
\item Polipropileno (PP).
\item Policloruro de vinilo (PVC).
\item Poliestireno (PS).
\item PET.
\item Nailon.
\item Teflón (PTFE).
\end{itemize}

Presentan una enorme importancia industrial debido a su bajo coste y facilidad de fabricación.

\end{enumerate}

%%%%%%%%%%%%%%%%%%%%%%%%%%%%%%%%%%%%%%%%%%%%%%%%%%%%%%%%%%%%%

\subsection{Clasificación según su composición}

Según el número de monómeros diferentes que participan en la polimerización, los polímeros pueden clasificarse en:

\begin{description}

\item[Homopolímeros]

Se forman por la repetición de un único tipo de monómero.

Por ejemplo:

\[
n\,CH_2=CH_2
\longrightarrow
(-CH_2-CH_2-)_n
\]

El polietileno es un homopolímero porque toda la cadena está formada por unidades de eteno.

\vspace{0.3cm}

\item[Heteropolímeros o copolímeros]

Se obtienen a partir de dos o más monómeros distintos.

Ejemplos:

\begin{itemize}

\item PET.

\item Nailon 6,6.

\item Resinas epoxi.

\end{itemize}

La combinación de diferentes monómeros permite modificar las propiedades del material final.

\end{description}

%%%%%%%%%%%%%%%%%%%%%%%%%%%%%%%%%%%%%%%%%%%%%%%%%%%%%%%%%%%%%

\subsection{Clasificación según su estructura}

La forma en que se organizan las cadenas poliméricas determina muchas de sus propiedades mecánicas.

\begin{enumerate}[label=\alph*)]

\item \textbf{Polímeros lineales}

Cada monómero se une únicamente a dos monómeros vecinos formando largas cadenas.

Presentan:

\begin{itemize}

\item Mayor flexibilidad.

\item Buena capacidad de deformación.

\item Pueden fundirse con facilidad.

\end{itemize}

Ejemplos:

PEAD, PVC y nailon.

\vspace{0.4cm}

\item \textbf{Polímeros ramificados}

Algunos monómeros se unen a tres o más posiciones formando ramas laterales.

Estas ramas dificultan el empaquetamiento de las cadenas y disminuyen la densidad del material.

Ejemplo:

PEBD.

\vspace{0.4cm}

\item \textbf{Polímeros reticulados}

Las cadenas se encuentran unidas mediante numerosos enlaces covalentes formando una red tridimensional.

Presentan:

\begin{itemize}

\item Gran rigidez.

\item Elevada resistencia térmica.

\item No pueden fundirse.

\end{itemize}

Ejemplos:

\begin{itemize}

\item Baquelita.

\item Resinas epoxi.

\end{itemize}

\end{enumerate}

%%%%%%%%%%%%%%%%%%%%%%%%%%%%%%%%%%%%%%%%%%%%%%%%%%%%%%%%%%%%%

\subsection{Clasificación según su comportamiento frente al calor}

Cuando aumenta la temperatura, los polímeros no responden todos de la misma forma.

\begin{center}
\begin{tabular}{p{3.5cm}p{7.5cm}}
\toprule
\textbf{Tipo} & \textbf{Características} \\
\midrule

Termoplásticos &
Se ablandan al calentarse y recuperan su forma al enfriarse. Pueden reciclarse mediante fusión.\\

\addlinespace

Termoestables &
Durante su fabricación forman enlaces entre cadenas que impiden volver a fundirlos. Son muy resistentes al calor.\\

\addlinespace

Elastómeros &
Presentan una ligera reticulación que les permite deformarse mucho y recuperar posteriormente su forma original.\\

\bottomrule
\end{tabular}
\end{center}

\vspace{0.3cm}

Ejemplos representativos:

\begin{itemize}

\item Termoplásticos: PE, PP, PVC, PET.

\item Termoestables: baquelita, melamina.

\item Elastómeros: caucho natural, neopreno y silicona.

\end{itemize}

%%%%%%%%%%%%%%%%%%%%%%%%%%%%%%%%%%%%%%%%%%%%%%%%%%%%%%%%%%%%%

\begin{tcolorbox}[
title=Recuerda,
colback=blue!5,
colframe=blue!70!black]

La clasificación de un polímero depende de varios criterios al mismo tiempo. Por ejemplo, el PET es un polímero sintético, copolímero, lineal y termoplástico.

\end{tcolorbox}

%%%%%%%%%%%%%%%%%%%%%%%%%%%%%%%%%%%%%%%%%%%%%%%%%%%%%%%%%%%%%

\begin{tcolorbox}[
title=Curiosidad,
colback=green!5,
colframe=green!50!black]

Aunque solemos asociar los polímeros con los plásticos, muchas biomoléculas esenciales para la vida también son polímeros. El ADN, las proteínas y la celulosa son algunos de los ejemplos más importantes.

\end{tcolorbox}

%
\begin{figure}[H]
\centering

\begin{tikzpicture}[
>=Latex,
font=\small,

main/.style={
rectangle,
rounded corners=3pt,
draw=DarkTeal,
very thick,
fill=DarkTeal!18,
minimum width=3.8cm,
minimum height=0.9cm,
align=center,
drop shadow
},

group/.style={
rectangle,
rounded corners=3pt,
draw=blue!70!black,
thick,
fill=blue!10,
minimum width=3.6cm,
minimum height=0.8cm,
align=center,
drop shadow
},

item/.style={
rectangle,
rounded corners=2pt,
draw=gray!70,
fill=gray!8,
minimum width=3.6cm,
minimum height=0.7cm,
align=left
},

line/.style={
-{Latex[length=3mm]},
thick
}
]

%%%%%%%%%%%%%%%%%%%%%%%%%%%%%%%%%%%%%%%%%%%%%%%%%%%%%%%%%%%%
% Nodo principal
%%%%%%%%%%%%%%%%%%%%%%%%%%%%%%%%%%%%%%%%%%%%%%%%%%%%%%%%%%%%

\node[main] (P) {POLÍMEROS};

%%%%%%%%%%%%%%%%%%%%%%%%%%%%%%%%%%%%%%%%%%%%%%%%%%%%%%%%%%%%
% Nivel 2
%%%%%%%%%%%%%%%%%%%%%%%%%%%%%%%%%%%%%%%%%%%%%%%%%%%%%%%%%%%%

\node[group,below left=2cm and 4.8cm of P] (O)
{Según su\\origen};

\node[group,below left=2cm and 1.5cm of P] (C)
{Según su\\composición};

\node[group,below right=2cm and 1.5cm of P] (E)
{Según su\\estructura};

\node[group,below right=2cm and 4.8cm of P] (T)
{Según su\\comportamiento térmico};

%%%%%%%%%%%%%%%%%%%%%%%%%%%%%%%%%%%%%%%%%%%%%%%%%%%%%%%%%%%%

\draw[line](P)--(O);
\draw[line](P)--(C);
\draw[line](P)--(E);
\draw[line](P)--(T);

%%%%%%%%%%%%%%%%%%%%%%%%%%%%%%%%%%%%%%%%%%%%%%%%%%%%%%%%%%%%
% ORIGEN
%%%%%%%%%%%%%%%%%%%%%%%%%%%%%%%%%%%%%%%%%%%%%%%%%%%%%%%%%%%%

\node[item,below=8mm of O]
(O1)
{• Naturales\\
Celulosa, proteínas, ADN, caucho};

\node[item,below=2mm of O1]
(O2)
{• Sintéticos\\
PE, PVC, PET, nailon};

\draw[line](O)--(O1);
\draw[line](O1)--(O2);

%%%%%%%%%%%%%%%%%%%%%%%%%%%%%%%%%%%%%%%%%%%%%%%%%%%%%%%%%%%%
% COMPOSICIÓN
%%%%%%%%%%%%%%%%%%%%%%%%%%%%%%%%%%%%%%%%%%%%%%%%%%%%%%%%%%%%

\node[item,below=8mm of C]
(C1)
{• Homopolímeros\\
Un único monómero};

\node[item,below=2mm of C1]
(C2)
{• Copolímeros\\
Dos o más monómeros};

\draw[line](C)--(C1);
\draw[line](C1)--(C2);

%%%%%%%%%%%%%%%%%%%%%%%%%%%%%%%%%%%%%%%%%%%%%%%%%%%%%%%%%%%%
% ESTRUCTURA
%%%%%%%%%%%%%%%%%%%%%%%%%%%%%%%%%%%%%%%%%%%%%%%%%%%%%%%%%%%%

\node[item,below=8mm of E]
(E1)
{• Lineales};

\node[item,below=2mm of E1]
(E2)
{• Ramificados};

\node[item,below=2mm of E2]
(E3)
{• Reticulados};

\draw[line](E)--(E1);
\draw[line](E1)--(E2);
\draw[line](E2)--(E3);

%%%%%%%%%%%%%%%%%%%%%%%%%%%%%%%%%%%%%%%%%%%%%%%%%%%%%%%%%%%%
% TÉRMICO
%%%%%%%%%%%%%%%%%%%%%%%%%%%%%%%%%%%%%%%%%%%%%%%%%%%%%%%%%%%%

\node[item,below=8mm of T]
(T1)
{• Termoplásticos};

\node[item,below=2mm of T1]
(T2)
{• Termoestables};

\node[item,below=2mm of T2]
(T3)
{• Elastómeros};

\draw[line](T)--(T1);
\draw[line](T1)--(T2);
\draw[line](T2)--(T3);

%%%%%%%%%%%%%%%%%%%%%%%%%%%%%%%%%%%%%%%%%%%%%%%%%%%%%%%%%%%%

\end{tikzpicture}

\caption{Principales criterios de clasificación de los polímeros.}

\end{figure}

%%%%%%%%%%%%%%%%%%%%%%%%%%%%%%%%%%%%%%%%%%%%%%%%%%%%%%%%%%%%
%%%%%%%%%%%%%%%%%%%%%%% SECCIÓN 3 %%%%%%%%%%%%%%%%%%%%%%%%%%%
%%%%%%%%%%%%%%%%%%%%%%%%%%%%%%%%%%%%%%%%%%%%%%%%%%%%%%%%%%%%

\section{Formación de los polímeros}

Los polímeros se obtienen mediante un conjunto de reacciones químicas denominadas
\textbf{reacciones de polimerización}. Durante estas reacciones, un gran número de
moléculas pequeñas (monómeros) se unen mediante enlaces covalentes para formar una
única macromolécula de gran tamaño.

En términos generales, el proceso puede representarse mediante la expresión:

\[
n\,\mathrm{M}\longrightarrow (\mathrm{M})_n
\]

donde:

\begin{itemize}
    \item \(M\) representa el monómero.
    \item \(n\) es el número de monómeros que forman la cadena.
\end{itemize}

Cuanto mayor sea el valor de \(n\), mayor será la longitud de la cadena y, en
general, mejores serán las propiedades mecánicas del polímero.

%%%%%%%%%%%%%%%%%%%%%%%%%%%%%%%%%%%%%%%%%%%%%%%%%%%%%%%%%%%%%

\subsection{Grado de polimerización}

El número de monómeros que forman una cadena recibe el nombre de
\textbf{grado de polimerización} (\(GP\)).

Matemáticamente se define como

\[
GP=\frac{M_{\rm polímero}}{M_{\rm monómero}}
\]

donde

\begin{itemize}
\item \(M_{\rm polímero}\) es la masa molecular media del polímero.
\item \(M_{\rm monómero}\) es la masa molecular del monómero.
\end{itemize}

En los polímeros industriales el grado de polimerización suele variar desde algunos
cientos hasta varios cientos de miles.

\begin{tcolorbox}[
title=Importante,
colback=yellow!8,
colframe=orange!70!black]

Dos muestras del mismo polímero pueden presentar propiedades diferentes si poseen
distintos grados de polimerización. Cuanto mayor es la longitud de las cadenas,
mayor suele ser la resistencia mecánica y la temperatura de reblandecimiento.

\end{tcolorbox}

%%%%%%%%%%%%%%%%%%%%%%%%%%%%%%%%%%%%%%%%%%%%%%%%%%%%%%%%%%%%%

\subsection{Esquema general de formación}

La formación de un polímero puede entenderse como una sucesión de uniones entre
monómeros.

\begin{center}

\begin{tikzpicture}[
    every node/.style={font=\small},
    box/.style={
        draw,
        rounded corners,
        minimum width=2.2cm,
        minimum height=8mm,
        align=center
    },
    node distance=6mm
]

\node[box] (M1) {Monómero};

\node[right=of M1] (plus) {$+$};

\node[box,right=of plus] (M2) {Monómero};

\node[right=of M2] (arrow) {$\Longrightarrow$};

\node[
    box,
    minimum width=3.8cm,
    right=of arrow
] (P) {$-\mathrm{M}-\mathrm{M}-\mathrm{M}-\cdots$};

\end{tikzpicture}

\end{center}

Durante la reacción, cada monómero aporta uno o varios enlaces que permiten su unión
con los monómeros vecinos, originando largas cadenas moleculares.

%%%%%%%%%%%%%%%%%%%%%%%%%%%%%%%%%%%%%%%%%%%%%%%%%%%%%%%%%%%%%

\subsection{Representación}

Un ejemplo sencillo es la formación del polietileno.

\[
\chemfig{CH_2=CH_2}
\qquad
\Longrightarrow
\qquad
\chemleft(
\chemfig{-CH_2-CH_2-}
\chemright)_n
\]

En esta reacción el doble enlace carbono-carbono desaparece y cada molécula de
eteno queda enlazada con las dos moléculas vecinas.

%%%%%%%%%%%%%%%%%%%%%%%%%%%%%%%%%%%%%%%%%%%%%%%%%%%%%%%%%%%%%

\subsection{Factores que influyen en la polimerización}

Las propiedades del polímero obtenido dependen de numerosos factores:

\begin{enumerate}[label=\alph*)]

\item Naturaleza química del monómero.

\item Temperatura de reacción.

\item Presión.

\item Catalizadores empleados.

\item Tiempo de reacción.

\item Presencia de impurezas.

\end{enumerate}

Modificando estas variables es posible obtener materiales con propiedades muy
diferentes a partir del mismo monómero.

%%%%%%%%%%%%%%%%%%%%%%%%%%%%%%%%%%%%%%%%%%%%%%%%%%%%%%%%%%%%%

\subsection{Relación entre estructura y propiedades}

Las propiedades de un polímero dependen no solamente del monómero utilizado, sino
también de cómo se organizan las cadenas.

\begin{center}

\begin{tabular}{p{4cm}p{8.2cm}}

\toprule

\textbf{Característica} & \textbf{Consecuencia} \\

\midrule

Cadenas largas &
Mayor resistencia mecánica. \\

\addlinespace

Mayor ramificación &
Disminuye la densidad y aumenta la flexibilidad. \\

\addlinespace

Mayor reticulación &
Incrementa la rigidez y la resistencia al calor. \\

\addlinespace

Mayor cristalinidad &
Aumenta la dureza y la temperatura de fusión. \\

\bottomrule

\end{tabular}

\end{center}

%%%%%%%%%%%%%%%%%%%%%%%%%%%%%%%%%%%%%%%%%%%%%%%%%%%%%%%%%%%%%

\subsection{Tipos de polimerización}

Dependiendo del mecanismo mediante el cual se unen los monómeros, las reacciones de
polimerización se clasifican en dos grandes grupos:

\begin{center}

\begin{tikzpicture}[
every node/.style={draw,rounded corners,align=center},
node distance=20mm]

\node (A){Polimerización};

\node[below left=of A,minimum width=4cm](B)
{Adición\\No se eliminan moléculas};

\node[below right=of A,minimum width=4cm](C)
{Condensación\\Se elimina una molécula pequeña};

\draw[-latex](A)--(B);

\draw[-latex](A)--(C);

\end{tikzpicture}

\end{center}

En los siguientes apartados se estudiarán con detalle ambos mecanismos de
polimerización y los polímeros más importantes que se obtienen mediante cada uno de
ellos.

%%%%%%%%%%%%%%%%%%%%%%%%%%%%%%%%%%%%%%%%%%%%%%%%%%%%%%%%%%%%%

\begin{tcolorbox}[
title=Recuerda,
colback=blue!5,
colframe=blue!70!black]

La polimerización consiste en unir un gran número de monómeros para formar una
macromolécula. Dependiendo de si durante el proceso se pierde o no una molécula
pequeña, distinguimos entre polimerización por adición y por condensación.

\end{tcolorbox}

%%%%%%%%%%%%%%%%%%%%%%%%%%%%%%%%%%%%%%%%%%%%%%%%%%%%%%%%%%%%
%%%%%%%%%%%%%%%%%%%%%% SECCIÓN 4A %%%%%%%%%%%%%%%%%%%%%%%%%%%
%%%%%%%%%%%%%%% POLIMERIZACIÓN POR ADICIÓN %%%%%%%%%%%%%%%%%%
%%%%%%%%%%%%%%%%%%%%%%%%%%%%%%%%%%%%%%%%%%%%%%%%%%%%%%%%%%%%

\section{Polimerización por adición}

La \textbf{polimerización por adición} es el proceso mediante el cual un gran número
de monómeros insaturados se unen para formar una macromolécula sin que se produzca la
pérdida de ningún átomo o molécula durante la reacción.

Los monómeros utilizados suelen contener un \textbf{doble} o un \textbf{triple enlace}
carbono-carbono. Durante la polimerización estos enlaces múltiples se rompen,
transformándose en enlaces sencillos que permiten la unión sucesiva de los monómeros.

En consecuencia,

\begin{itemize}
\item la masa del polímero es exactamente la suma de las masas de todos los
monómeros que lo forman;
\item no se generan productos secundarios;
\item el polímero conserva la misma composición elemental que el monómero.
\end{itemize}

Este tipo de polimerización es el más utilizado en la fabricación de plásticos de
gran consumo, como el polietileno o el PVC.

%%%%%%%%%%%%%%%%%%%%%%%%%%%%%%%%%%%%%%%%%%%%%%%%%%%%%%%%%%%%

\subsection{Características}

Las principales características de la polimerización por adición son:

\begin{enumerate}[label=\alph*)]

\item Los monómeros presentan enlaces múltiples (\(\mathrm{C=C}\) o
\(\mathrm{C\equiv C}\)).

\item No se desprenden moléculas pequeñas como agua, amoníaco o cloruro de
hidrógeno.

\item Se obtienen generalmente \textbf{homopolímeros}, ya que la cadena está formada
por la repetición del mismo monómero.

\item El proceso suele producir cadenas muy largas con masas moleculares muy
elevadas.

\end{enumerate}

%%%%%%%%%%%%%%%%%%%%%%%%%%%%%%%%%%%%%%%%%%%%%%%%%%%%%%%%%%%%

\subsection{Representación general}

La reacción general puede escribirse como

\[
n\,CH_2=CH-R
\longrightarrow
(-CH_2-CH(R)-)_n
\]

donde \(R\) representa cualquier átomo o grupo de átomos unido al doble enlace.

Cuando \(R=H\) se obtiene el polietileno; cuando \(R=Cl\), el PVC; y cuando
\(R=C_6H_5\), el poliestireno.

%%%%%%%%%%%%%%%%%%%%%%%%%%%%%%%%%%%%%%%%%%%%%%%%%%%%%%%%%%%%

\subsection{Mecanismo radicalario}

La mayoría de las polimerizaciones por adición tienen lugar mediante un mecanismo
radicalario.

En él intervienen especies químicas muy reactivas denominadas
\textbf{radicales libres}, que poseen un electrón desapareado.

El proceso consta de tres etapas.

%%%%%%%%%%%%%%%%%%%%%%%%%%%%%%%%%%%%%%%%%%%%%%%%%%%%%%%%%%%%

\subsubsection*{1. Iniciación}

La reacción comienza gracias a un iniciador, normalmente un peróxido, que se rompe
por acción del calor o de la luz formando radicales libres.

\[
ROOR
\longrightarrow
2RO\cdot
\]

Estos radicales atacan el doble enlace del primer monómero, originando un nuevo
radical en el extremo de la cadena.

\vspace{0.4cm}

\begin{center}


\begin{tikzpicture}[
    node distance=40mm,
    every node/.style={
        draw,
        rounded corners,
        minimum width=3cm,
        minimum height=8mm
    },
    >={Latex}
]

\node (A) {Iniciador};
\node[right=of A] (C) {Radical libre};

\draw[-Latex,thick] (A.east) -- (C.west);

\end{tikzpicture}


\end{center}

%%%%%%%%%%%%%%%%%%%%%%%%%%%%%%%%%%%%%%%%%%%%%%%%%%%%%%%%%%%%

\subsubsection*{2. Propagación}

El radical formado reacciona con otro monómero, que pasa a formar parte de la cadena.

El nuevo extremo continúa siendo un radical libre, por lo que el proceso puede
repetirse miles de veces.

\[
\cdot CH_2-CH_2
+
CH_2=CH_2
\longrightarrow
CH_2-CH_2-CH_2-CH_2\cdot
\]

Cada nueva incorporación aumenta la longitud de la cadena.

Esta etapa es la responsable del gran tamaño que alcanzan los polímeros.

%%%%%%%%%%%%%%%%%%%%%%%%%%%%%%%%%%%%%%%%%%%%%%%%%%%%%%%%%%%%

\subsubsection*{3. Terminación}

El crecimiento de la cadena finaliza cuando desaparece el radical libre del extremo.

Esto puede suceder de varias formas:

\begin{itemize}

\item Unión de dos cadenas radicalarias.

\item Reacción con una impureza.

\item Transferencia de un átomo de hidrógeno.

\end{itemize}

Una vez terminada la reacción, el polímero ya no puede seguir creciendo.

%%%%%%%%%%%%%%%%%%%%%%%%%%%%%%%%%%%%%%%%%%%%%%%%%%%%%%%%%%%%

\subsection{Esquema del mecanismo}

\begin{center}

\begin{tikzpicture}[node distance=12mm]

\node[draw,rounded corners,minimum width=3cm](A){Iniciación};

\node[right=of A,draw,rounded corners,minimum width=3cm](B){Propagación};

\node[right=of B,draw,rounded corners,minimum width=3cm](C){Terminación};

\draw[-latex,thick](A)--(B);

\draw[-latex,thick](B)--(C);

\end{tikzpicture}

\end{center}

%%%%%%%%%%%%%%%%%%%%%%%%%%%%%%%%%%%%%%%%%%%%%%%%%%%%%%%%%%%%

\subsection{Ejemplo: formación del polietileno}

El polietileno se obtiene por polimerización del eteno.

\[
n\,CH_2=CH_2
\longrightarrow
(-CH_2-CH_2-)_n
\]

Durante la reacción desaparece el doble enlace y cada molécula de eteno queda unida
a las moléculas vecinas mediante enlaces sencillos.

El resultado es una cadena muy larga formada únicamente por átomos de carbono e
hidrógeno.

%%%%%%%%%%%%%%%%%%%%%%%%%%%%%%%%%%%%%%%%%%%%%%%%%%%%%%%%%%%%

\begin{tcolorbox}[
title=Recuerda,
colback=blue!5,
colframe=blue!70!black]

En una polimerización por adición:

\begin{itemize}

\item No se pierde ninguna molécula.

\item Los monómeros poseen dobles o triples enlaces.

\item La masa del polímero es igual a la suma de las masas de los monómeros.

\item Generalmente se obtienen homopolímeros.

\end{itemize}

\end{tcolorbox}

%%%%%%%%%%%%%%%%%%%%%%%%%%%%%%%%%%%%%%%%%%%%%%%%%%%%%%%%%%%%

\begin{tcolorbox}[
title=Curiosidad,
colback=green!5,
colframe=green!60!black]

El polietileno es el plástico más producido del mundo. Cada año se fabrican más de
100 millones de toneladas y se utiliza en bolsas, botellas, envases, tuberías y una
gran variedad de productos cotidianos.

\end{tcolorbox}

%%%%%%%%%%%%%%%%%%%%%%%%%%%%%%%%%%%%%%%%%%%%%%%%%%%%%%%%%%%%
%%%%%%%%%%%%%%%%%%%%%%% SECCIÓN 4B %%%%%%%%%%%%%%%%%%%%%%%%%%
%%%%%%%%%%%% PRINCIPALES POLÍMEROS DE ADICIÓN %%%%%%%%%%%%%%%
%%%%%%%%%%%%%%%%%%%%%%%%%%%%%%%%%%%%%%%%%%%%%%%%%%%%%%%%%%%%

\section{Principales polímeros de adición}

La polimerización por adición permite obtener algunos de los materiales más utilizados
en la actualidad. La mayoría de ellos son termoplásticos, ligeros, resistentes a la
corrosión y fáciles de moldear, lo que explica su enorme importancia industrial.

%%%%%%%%%%%%%%%%%%%%%%%%%%%%%%%%%%%%%%%%%%%%%%%%%%%%%%%%%%%%

\subsection{Polietileno (PE)}

El polietileno es el polímero de mayor producción mundial. Se obtiene por
polimerización del eteno.

\[
n\,CH_2=CH_2
\longrightarrow
(-CH_2-CH_2-)_n
\]

Dependiendo de las condiciones de fabricación pueden obtenerse dos tipos principales:

\begin{itemize}
\item \textbf{PEBD (polietileno de baja densidad):} presenta cadenas muy ramificadas,
es flexible y ligero.
\item \textbf{PEAD (polietileno de alta densidad):} posee cadenas casi lineales,
presentando mayor resistencia mecánica y rigidez.
\end{itemize}

\textbf{Aplicaciones}

\begin{itemize}
\item Bolsas de plástico.
\item Botellas.
\item Envases alimentarios.
\item Tuberías.
\item Aislamiento de cables.
\end{itemize}

%%%%%%%%%%%%%%%%%%%%%%%%%%%%%%%%%%%%%%%%%%%%%%%%%%%%%%%%%%%%

\subsection{Polipropileno (PP)}

Se obtiene mediante la polimerización del propeno.

\[
n\,CH_2=CH-CH_3
\longrightarrow
(-CH_2-CH(CH_3)-)_n
\]

Es un plástico muy ligero, rígido y resistente a numerosos productos químicos.

\textbf{Aplicaciones}

\begin{itemize}
\item Tapones.
\item Envases reutilizables.
\item Material sanitario.
\item Componentes de automóviles.
\item Mobiliario.
\end{itemize}

%%%%%%%%%%%%%%%%%%%%%%%%%%%%%%%%%%%%%%%%%%%%%%%%%%%%%%%%%%%%

\subsection{Policloruro de vinilo (PVC)}

Se obtiene a partir del cloroeteno (cloruro de vinilo).

\[
n\,CH_2=CHCl
\longrightarrow
(-CH_2-CHCl-)_n
\]

La presencia del átomo de cloro aumenta su dureza y resistencia al envejecimiento.

Según los aditivos empleados puede obtenerse un PVC rígido o flexible.

\textbf{Aplicaciones}

\begin{itemize}
\item Tuberías.
\item Ventanas.
\item Aislamiento eléctrico.
\item Tarjetas.
\item Suelos vinílicos.
\end{itemize}

%%%%%%%%%%%%%%%%%%%%%%%%%%%%%%%%%%%%%%%%%%%%%%%%%%%%%%%%%%%%

\subsection{Poliestireno (PS)}

Procede de la polimerización del estireno.

\[
n\,CH_2=CH(C_6H_5)
\longrightarrow
(-CH_2-CH(C_6H_5)-)_n
\]

Es un material rígido y transparente.

Cuando contiene gran cantidad de aire recibe el nombre de
\textbf{poliestireno expandido (EPS)}, conocido comercialmente como
\emph{corcho blanco}.

\textbf{Aplicaciones}

\begin{itemize}
\item Embalajes.
\item Vasos desechables.
\item Bandejas para alimentos.
\item Aislamiento térmico.
\end{itemize}

%%%%%%%%%%%%%%%%%%%%%%%%%%%%%%%%%%%%%%%%%%%%%%%%%%%%%%%%%%%%

\subsection{Politetrafluoroetileno (PTFE)}

El PTFE, conocido comercialmente como \textbf{Teflón}, se obtiene a partir del
tetrafluoroeteno.

\[
n\,CF_2=CF_2
\longrightarrow
(-CF_2-CF_2-)_n
\]

Posee una extraordinaria resistencia química y un coeficiente de rozamiento muy bajo.

\textbf{Aplicaciones}

\begin{itemize}
\item Recubrimientos antiadherentes.
\item Juntas.
\item Válvulas.
\item Material de laboratorio.
\item Industria química.
\end{itemize}

%%%%%%%%%%%%%%%%%%%%%%%%%%%%%%%%%%%%%%%%%%%%%%%%%%%%%%%%%%%%

\subsection{Comparación de los principales polímeros}

\begin{center}

\renewcommand{\arraystretch}{1.3}

\begin{tabular}{p{2.6cm}p{2.9cm}p{3cm}p{4cm}}

\toprule

\textbf{Polímero} &
\textbf{Monómero} &
\textbf{Propiedad destacada} &
\textbf{Aplicaciones} \\

\midrule

Polietileno (PE)

& Eteno

& Flexible y ligero

& Bolsas, botellas, envases \\

\addlinespace

Polipropileno (PP)

& Propeno

& Muy resistente

& Tapones, recipientes, mobiliario \\

\addlinespace

PVC

& Cloroeteno

& Muy resistente al desgaste

& Tuberías, ventanas, cables \\

\addlinespace

Poliestireno (PS)

& Estireno

& Ligero y aislante

& Embalajes y envases \\

\addlinespace

PTFE

& Tetrafluoroeteno

& Antiadherente y resistente

& Sartenes, válvulas y juntas \\

\bottomrule

\end{tabular}

\end{center}

%%%%%%%%%%%%%%%%%%%%%%%%%%%%%%%%%%%%%%%%%%%%%%%%%%%%%%%%%%%%

\subsection{Ventajas de los polímeros de adición}

La enorme difusión de estos materiales se debe a que presentan numerosas ventajas:

\begin{itemize}

\item Baja densidad.

\item Elevada resistencia química.

\item Buenos aislantes eléctricos.

\item Facilidad de fabricación.

\item Bajo coste.

\item Gran variedad de propiedades.

\end{itemize}

No obstante, muchos de ellos son poco biodegradables, por lo que resulta imprescindible
su correcta reutilización y reciclaje.

%%%%%%%%%%%%%%%%%%%%%%%%%%%%%%%%%%%%%%%%%%%%%%%%%%%%%%%%%%%%

\begin{tcolorbox}[
title=Importante,
colback=blue!5,
colframe=blue!70!black]

Todos los polímeros estudiados en esta sección se obtienen mediante
\textbf{polimerización por adición}. En ellos no se elimina ninguna molécula durante
la reacción y la unidad repetitiva conserva todos los átomos del monómero original.

\end{tcolorbox}

%%%%%%%%%%%%%%%%%%%%%%%%%%%%%%%%%%%%%%%%%%%%%%%%%%%%%%%%%%%%

\begin{tcolorbox}[
title=Curiosidad,
colback=green!5,
colframe=green!60!black]

Aunque el Teflón es conocido por su uso en utensilios de cocina, también se emplea en
la industria aeroespacial, en equipos médicos y en componentes electrónicos debido a
su extraordinaria estabilidad química y térmica.

\end{tcolorbox}

%%%%%%%%%%%%%%%%%%%%%%%%%%%%%%%%%%%%%%%%%%%%%%%%%%%%%%%%%%%%
%%%%%%%%%%%%%%%%%%%%%%%% SECCIÓN 5 %%%%%%%%%%%%%%%%%%%%%%%%%%
%%%%%%%%%%%%%% POLIMERIZACIÓN POR CONDENSACIÓN %%%%%%%%%%%%%%
%%%%%%%%%%%%%%%%%%%%%%%%%%%%%%%%%%%%%%%%%%%%%%%%%%%%%%%%%%%%

\section{Polimerización por condensación}

La \textbf{polimerización por condensación} consiste en la unión sucesiva de dos o más
tipos de monómeros mediante reacciones en las que, además del polímero, se forma una
molécula pequeña como agua (\(\mathrm{H_2O}\)), cloruro de hidrógeno
(\(\mathrm{HCl}\)) o metanol (\(\mathrm{CH_3OH}\)).

A diferencia de la polimerización por adición, en este proceso la masa del polímero
\textbf{no coincide} exactamente con la suma de las masas de los monómeros, ya que
parte de la materia se elimina en forma de moléculas pequeñas.

En la mayoría de los casos intervienen monómeros que poseen dos grupos funcionales,
lo que permite que la reacción continúe en ambos extremos de la cadena.

%%%%%%%%%%%%%%%%%%%%%%%%%%%%%%%%%%%%%%%%%%%%%%%%%%%%%%%%%%%%

\subsection{Características}

Los polímeros obtenidos por condensación presentan las siguientes características:

\begin{itemize}

\item Se forman mediante reacciones entre grupos funcionales.

\item Durante la reacción se elimina una molécula pequeña.

\item Habitualmente intervienen dos monómeros diferentes.

\item Suelen originar \textbf{copolímeros}.

\item Muchos presentan una elevada resistencia mecánica y térmica.

\end{itemize}

%%%%%%%%%%%%%%%%%%%%%%%%%%%%%%%%%%%%%%%%%%%%%%%%%%%%%%%%%%%%

\subsection{Comparación entre adición y condensación}

\begin{center}

\renewcommand{\arraystretch}{1.3}

\begin{tabular}{p{5cm}p{5cm}}

\toprule

\textbf{Polimerización por adición} &
\textbf{Polimerización por condensación}

\\

\midrule

No se elimina ninguna molécula.

&
Se elimina una molécula pequeña.

\\

Monómeros con dobles enlaces.

&
Monómeros con grupos funcionales.

\\

Generalmente homopolímeros.

&
Generalmente copolímeros.

\\

La masa del polímero coincide con la suma de los monómeros.

&
La masa disminuye debido a la molécula eliminada.

\\

Ejemplos: PE, PVC, PP, PS.

&
Ejemplos: PET, nailon, baquelita.

\\

\bottomrule

\end{tabular}

\end{center}

%%%%%%%%%%%%%%%%%%%%%%%%%%%%%%%%%%%%%%%%%%%%%%%%%%%%%%%%%%%%

\subsection{Formación de un poliéster}

Los \textbf{poliésteres} se obtienen por reacción entre un diácido y un diol.

Cada vez que se forma un enlace éster se libera una molécula de agua.

De forma general:

\[
n\,HO-R-OH
+
n\,HOOC-R'-COOH
\longrightarrow
(-O-R-O-CO-R'-CO-)_n
+
2nH_2O
\]

%%%%%%%%%%%%%%%%%%%%%%%%%%%%%%%%%%%%%%%%%%%%%%%%%%%%%%%%%%%%

\subsubsection*{El PET}

El poli(etileno tereftalato), conocido como \textbf{PET}, es el poliéster más
importante desde el punto de vista industrial.

Se obtiene mediante la reacción entre:

\begin{itemize}

\item ácido tereftálico;

\item etilenglicol (etano-1,2-diol).

\end{itemize}

Sus principales propiedades son:

\begin{itemize}

\item Transparencia.

\item Gran resistencia mecánica.

\item Buena resistencia química.

\item Baja permeabilidad a los gases.

\end{itemize}

\textbf{Aplicaciones}

\begin{itemize}

\item Botellas de bebidas.

\item Envases alimentarios.

\item Fibras textiles (tergal).

\item Láminas plásticas.

\end{itemize}

%%%%%%%%%%%%%%%%%%%%%%%%%%%%%%%%%%%%%%%%%%%%%%%%%%%%%%%%%%%%

\subsection{Formación de una poliamida}

Las \textbf{poliamidas} contienen enlaces amida
(\(-CONH-\)) en la cadena principal.

Se obtienen por reacción entre un diácido y una diamina.

Cada nuevo enlace amida genera una molécula de agua.

De forma simplificada:

\[
n\,H_2N-R-NH_2
+
n\,HOOC-R'-COOH
\longrightarrow
(-NH-R-NH-CO-R'-CO-)_n
+
2nH_2O
\]

%%%%%%%%%%%%%%%%%%%%%%%%%%%%%%%%%%%%%%%%%%%%%%%%%%%%%%%%%%%%

\subsubsection*{Nailon 6,6}

El nailon 6,6 es la poliamida más conocida.

Se obtiene a partir de:

\begin{itemize}

\item ácido adípico;

\item hexano-1,6-diamina.

\end{itemize}

Presenta:

\begin{itemize}

\item elevada resistencia mecánica;

\item gran elasticidad;

\item resistencia al desgaste;

\item bajo rozamiento.

\end{itemize}

\textbf{Aplicaciones}

\begin{itemize}

\item Tejidos.

\item Cuerdas.

\item Sedales.

\item Paracaídas.

\item Engranajes.

\item Cepillos.

\end{itemize}

%%%%%%%%%%%%%%%%%%%%%%%%%%%%%%%%%%%%%%%%%%%%%%%%%%%%%%%%%%%%

\subsection{Otros polímeros de condensación}

Además del PET y del nailon, existen otros polímeros muy utilizados.

\subsubsection*{Baquelita}

Fue el primer plástico completamente sintético.

Es un polímero termoestable obtenido por condensación del fenol y el formaldehído.

Se emplea en:

\begin{itemize}

\item enchufes;

\item interruptores;

\item mangos de utensilios;

\item componentes eléctricos.

\end{itemize}

\vspace{0.3cm}

\subsubsection*{Poliuretanos}

Se obtienen por reacción entre polioles e isocianatos.

Dependiendo de su estructura pueden fabricarse:

\begin{itemize}

\item espumas flexibles;

\item espumas rígidas;

\item adhesivos;

\item pinturas;

\item aislantes térmicos.

\end{itemize}

%%%%%%%%%%%%%%%%%%%%%%%%%%%%%%%%%%%%%%%%%%%%%%%%%%%%%%%%%%%%

\subsection{Importancia industrial}

Los polímeros de condensación poseen aplicaciones muy variadas debido a sus buenas
propiedades mecánicas y químicas.

Muchas fibras textiles modernas, materiales compuestos y envases alimentarios están
fabricados con este tipo de polímeros.

Además, algunos de ellos pueden reciclarse mediante procesos químicos que permiten
recuperar parcialmente los monómeros originales.

%%%%%%%%%%%%%%%%%%%%%%%%%%%%%%%%%%%%%%%%%%%%%%%%%%%%%%%%%%%%

\begin{tcolorbox}[
title=Recuerda,
colback=blue!5,
colframe=blue!70!black]

En la polimerización por condensación siempre se produce la eliminación de una
molécula pequeña (normalmente agua). Los monómeros poseen grupos funcionales
reactivos y, generalmente, el polímero obtenido es un copolímero.

\end{tcolorbox}

%%%%%%%%%%%%%%%%%%%%%%%%%%%%%%%%%%%%%%%%%%%%%%%%%%%%%%%%%%%%

\begin{tcolorbox}[
title=Curiosidad,
colback=green!5,
colframe=green!60!black]

Las fibras de nailon revolucionaron la industria textil durante el siglo XX por su
gran resistencia y ligereza. Actualmente siguen empleándose en ropa deportiva,
paracaídas, airbags y numerosos componentes mecánicos.

\end{tcolorbox}

%%%%%%%%%%%%%%%%%%%%%%%%%%%%%%%%%%%%%%%%%%%%%%%%%%%%%%%%%%%%
%%%%%%%%%%%%%%%%%%%%%%%% SECCIÓN 6 %%%%%%%%%%%%%%%%%%%%%%%%%%
%%%%%%%%%%%%%%%%%% PROPIEDADES DE LOS POLÍMEROS %%%%%%%%%%%%%
%%%%%%%%%%%%%%%%%%%%%%%%%%%%%%%%%%%%%%%%%%%%%%%%%%%%%%%%%%%%

\section{Propiedades de los polímeros}

Los polímeros presentan una gran variedad de propiedades físicas y químicas que
dependen de su composición, de la longitud de sus cadenas y de las fuerzas que
actúan entre ellas. Estas propiedades determinan sus aplicaciones industriales y
su comportamiento frente al calor, la deformación o los agentes químicos.

En general, los polímeros poseen una menor densidad que los metales y las
cerámicas, son buenos aislantes y pueden fabricarse con propiedades muy diversas.

%%%%%%%%%%%%%%%%%%%%%%%%%%%%%%%%%%%%%%%%%%%%%%%%%%%%%%%%%%%%

\subsection{Factores que determinan sus propiedades}

Las propiedades de un polímero dependen principalmente de los siguientes factores.

\subsubsection*{Naturaleza del monómero}

La composición química del monómero determina los grupos funcionales presentes en
la cadena y, por tanto, las interacciones intermoleculares.

Por ejemplo:

\begin{itemize}

\item El polietileno, formado únicamente por carbono e hidrógeno, es flexible y
ligero.

\item El PVC contiene átomos de cloro, lo que aumenta su dureza y resistencia.

\item Las poliamidas presentan enlaces amida capaces de formar puentes de
hidrógeno, aumentando considerablemente su resistencia mecánica.

\end{itemize}

%%%%%%%%%%%%%%%%%%%%%%%%%%%%%%%%%%%%%%%%%%%%%%%%%%%%%%%%%%%%

\subsubsection*{Longitud de la cadena}

En general, cuanto mayor es el grado de polimerización:

\begin{itemize}

\item aumenta la resistencia mecánica;

\item aumenta la viscosidad cuando el material está fundido;

\item aumenta la temperatura de reblandecimiento.

\end{itemize}

%%%%%%%%%%%%%%%%%%%%%%%%%%%%%%%%%%%%%%%%%%%%%%%%%%%%%%%%%%%%

\subsubsection*{Ramificación}

La presencia de ramas dificulta el empaquetamiento de las cadenas.

Como consecuencia:

\begin{itemize}

\item disminuye la densidad;

\item aumenta la flexibilidad;

\item disminuye la resistencia mecánica.

\end{itemize}

Por ello, el polietileno de baja densidad (PEBD) es más flexible que el
polietileno de alta densidad (PEAD).

%%%%%%%%%%%%%%%%%%%%%%%%%%%%%%%%%%%%%%%%%%%%%%%%%%%%%%%%%%%%

\subsubsection*{Reticulación}

Cuando existen numerosos enlaces covalentes entre cadenas diferentes se forma una
estructura tridimensional denominada \textbf{red reticulada}.

Este tipo de estructura produce:

\begin{itemize}

\item gran rigidez;

\item elevada resistencia al calor;

\item imposibilidad de volver a fundir el material.

\end{itemize}

Los polímeros termoestables presentan un elevado grado de reticulación.

%%%%%%%%%%%%%%%%%%%%%%%%%%%%%%%%%%%%%%%%%%%%%%%%%%%%%%%%%%%%

\subsection{Propiedades físicas}

Las propiedades físicas más importantes son:

\begin{enumerate}[label=\alph*)]

\item \textbf{Baja densidad}

La mayoría de los polímeros presentan densidades comprendidas entre
\(0.9\) y \(1.5\ \mathrm{g\,cm^{-3}}\), muy inferiores a las de los metales.

\vspace{0.2cm}

\item \textbf{Elasticidad}

Muchos polímeros pueden deformarse considerablemente antes de romperse.

Los elastómeros son el mejor ejemplo de esta propiedad.

\vspace{0.2cm}

\item \textbf{Flexibilidad}

Las cadenas pueden deslizar unas sobre otras cuando las fuerzas intermoleculares
son débiles.

\vspace{0.2cm}

\item \textbf{Transparencia}

Algunos polímeros, como el PET o el policarbonato, son transparentes y permiten
el paso de la luz.

\vspace{0.2cm}

\item \textbf{Baja conductividad térmica}

Los polímeros son excelentes aislantes térmicos, por lo que se utilizan en la
fabricación de materiales aislantes.

\vspace{0.2cm}

\item \textbf{Baja conductividad eléctrica}

La mayoría son buenos aislantes eléctricos y se utilizan para recubrir cables y
fabricar componentes electrónicos.

\end{enumerate}

%%%%%%%%%%%%%%%%%%%%%%%%%%%%%%%%%%%%%%%%%%%%%%%%%%%%%%%%%%%%

\subsection{Propiedades químicas}

Desde el punto de vista químico destacan las siguientes características:

\begin{itemize}

\item Buena resistencia frente a numerosos ácidos y bases.

\item Resistencia a la corrosión.

\item Baja reactividad química.

\item Algunos polímeros pueden degradarse por acción de la radiación
ultravioleta o del oxígeno atmosférico.

\end{itemize}

Para aumentar su durabilidad suelen añadirse estabilizantes y antioxidantes
durante su fabricación.

%%%%%%%%%%%%%%%%%%%%%%%%%%%%%%%%%%%%%%%%%%%%%%%%%%%%%%%%%%%%

\subsection{Comportamiento frente al calor}

Según su respuesta al aumento de temperatura, los polímeros pueden clasificarse
en tres grandes grupos.

\begin{center}

\renewcommand{\arraystretch}{1.4}

\begin{tabular}{p{3.2cm}p{8.8cm}}

\toprule

\textbf{Tipo} &
\textbf{Características}

\\

\midrule

Termoplásticos

& Se ablandan al calentarse y pueden moldearse repetidamente.

\\

Termoestables

& Forman una red tridimensional y no vuelven a fundirse.

\\

Elastómeros

& Se deforman fácilmente y recuperan su forma inicial al cesar la fuerza.

\\

\bottomrule

\end{tabular}

\end{center}

%%%%%%%%%%%%%%%%%%%%%%%%%%%%%%%%%%%%%%%%%%%%%%%%%%%%%%%%%%%%

\subsection{Ventajas e inconvenientes}

\begin{minipage}{0.48\textwidth}

\textbf{Ventajas}

\begin{itemize}

\item Baja densidad.

\item Bajo coste.

\item Fácil fabricación.

\item Resistencia química.

\item Buenos aislantes.

\item Gran variedad de aplicaciones.

\end{itemize}

\end{minipage}
\hfill
\begin{minipage}{0.48\textwidth}

\textbf{Inconvenientes}

\begin{itemize}

\item Escasa resistencia a altas temperaturas.

\item Muchos proceden del petróleo.

\item Difícil degradación.

\item Formación de microplásticos.

\item Contaminación ambiental.

\end{itemize}

\end{minipage}

%%%%%%%%%%%%%%%%%%%%%%%%%%%%%%%%%%%%%%%%%%%%%%%%%%%%%%%%%%%%

\subsection{Relación entre estructura y propiedades}

La siguiente tabla resume cómo influye la estructura molecular en el
comportamiento del polímero.

\begin{center}

\renewcommand{\arraystretch}{1.3}

\begin{tabular}{p{5cm}p{7cm}}

\toprule

\textbf{Característica estructural} &
\textbf{Consecuencia}

\\

\midrule

Cadenas largas

& Mayor resistencia mecánica.

\\

Mayor ramificación

& Más flexibilidad y menor densidad.

\\

Mayor reticulación

& Mayor rigidez y resistencia térmica.

\\

Puentes de hidrógeno

& Mayor dureza y temperatura de fusión.

\\

Mayor cristalinidad

& Mayor dureza y resistencia.

\\

\bottomrule

\end{tabular}

\end{center}

%%%%%%%%%%%%%%%%%%%%%%%%%%%%%%%%%%%%%%%%%%%%%%%%%%%%%%%%%%%%

\begin{tcolorbox}[
title=Recuerda,
colback=blue!5,
colframe=blue!70!black]

Las propiedades de un polímero no dependen únicamente del monómero empleado.
También influyen la longitud de las cadenas, su grado de ramificación, la
reticulación y las fuerzas intermoleculares existentes entre ellas.

\end{tcolorbox}

%%%%%%%%%%%%%%%%%%%%%%%%%%%%%%%%%%%%%%%%%%%%%%%%%%%%%%%%%%%%

\begin{tcolorbox}[
title=Ejemplo,
colback=green!5,
colframe=green!60!black]

El PEAD y el PEBD están formados por el mismo monómero (eteno). Sin embargo,
el PEAD posee cadenas casi lineales, mientras que el PEBD presenta numerosas
ramificaciones. Esta diferencia estructural explica que el PEAD sea más rígido
y resistente, mientras que el PEBD es más flexible.

\end{tcolorbox}

%%%%%%%%%%%%%%%%%%%%%%%%%%%%%%%%%%%%%%%%%%%%%%%%%%%%%%%%%%%%
%%%%%%%%%%%%%%%%%%%%%%%% SECCIÓN 7 %%%%%%%%%%%%%%%%%%%%%%%%%%
%%%%%%%%%%% APLICACIONES, RECICLAJE E IMPACTO AMBIENTAL %%%%
%%%%%%%%%%%%%%%%%%%%%%%%%%%%%%%%%%%%%%%%%%%%%%%%%%%%%%%%%%%%

\section{Aplicaciones, impacto medioambiental y reciclaje}

Los polímeros forman parte de prácticamente todos los objetos que utilizamos a diario.
Su combinación de ligereza, resistencia, facilidad de fabricación y bajo coste ha
hecho posible el desarrollo de numerosos avances tecnológicos. Sin embargo, su uso
masivo también ha generado importantes problemas medioambientales.

%%%%%%%%%%%%%%%%%%%%%%%%%%%%%%%%%%%%%%%%%%%%%%%%%%%%%%%%%%%%

\subsection{Aplicaciones de los polímeros}

Actualmente los polímeros están presentes en casi todos los sectores industriales.

\subsubsection*{Industria del envase}

Los plásticos permiten fabricar envases ligeros, resistentes y económicos.

Ejemplos:

\begin{itemize}

\item Botellas de PET.

\item Bolsas de polietileno.

\item Envases de polipropileno.

\item Bandejas de poliestireno expandido.

\end{itemize}

%%%%%%%%%%%%%%%%%%%%%%%%%%%%%%%%%%%%%%%%%%%%%%%%%%%%%%%%%%%%

\subsubsection*{Construcción}

En construcción se emplean polímeros debido a su resistencia frente a la corrosión
y al agua.

Algunas aplicaciones son:

\begin{itemize}

\item Tuberías de PVC.

\item Aislantes térmicos.

\item Pinturas.

\item Espumas de poliuretano.

\item Ventanas y perfiles.

\end{itemize}

%%%%%%%%%%%%%%%%%%%%%%%%%%%%%%%%%%%%%%%%%%%%%%%%%%%%%%%%%%%%

\subsubsection*{Industria textil}

Muchas fibras textiles actuales son materiales poliméricos.

Entre las más utilizadas destacan:

\begin{itemize}

\item Nailon.

\item Poliéster.

\item Elastano.

\item Acrílicos.

\end{itemize}

Estas fibras presentan gran resistencia, elasticidad y facilidad de lavado.

%%%%%%%%%%%%%%%%%%%%%%%%%%%%%%%%%%%%%%%%%%%%%%%%%%%%%%%%%%%%

\subsubsection*{Medicina}

Los polímeros han revolucionado numerosos tratamientos médicos.

Se emplean en:

\begin{itemize}

\item Prótesis.

\item Suturas.

\item Catéteres.

\item Jeringuillas.

\item Bolsas para suero.

\item Impresión 3D de tejidos.

\end{itemize}

%%%%%%%%%%%%%%%%%%%%%%%%%%%%%%%%%%%%%%%%%%%%%%%%%%%%%%%%%%%%

\subsubsection*{Automoción y aeronáutica}

Muchos componentes metálicos han sido sustituidos por polímeros debido a su menor
peso.

Esto permite reducir el consumo de combustible y las emisiones de dióxido de carbono.

%%%%%%%%%%%%%%%%%%%%%%%%%%%%%%%%%%%%%%%%%%%%%%%%%%%%%%%%%%%%

\subsection{Ventajas del uso de polímeros}

El éxito industrial de estos materiales se debe a sus numerosas ventajas.

\begin{itemize}

\item Baja densidad.

\item Elevada resistencia química.

\item Buenos aislantes térmicos y eléctricos.

\item Gran facilidad de fabricación.

\item Bajo coste.

\item Amplia variedad de propiedades.

\end{itemize}

%%%%%%%%%%%%%%%%%%%%%%%%%%%%%%%%%%%%%%%%%%%%%%%%%%%%%%%%%%%%

\subsection{Problemas medioambientales}

El enorme consumo de plásticos también genera importantes impactos sobre el medio
ambiente.

Los principales problemas son:

\begin{enumerate}

\item Consumo de recursos fósiles para fabricar muchos polímeros.

\item Acumulación de residuos durante cientos de años.

\item Formación de microplásticos.

\item Contaminación de océanos, ríos y suelos.

\item Ingestión accidental por numerosos animales.

\item Emisión de gases contaminantes cuando se queman de forma incontrolada.

\end{enumerate}

%%%%%%%%%%%%%%%%%%%%%%%%%%%%%%%%%%%%%%%%%%%%%%%%%%%%%%%%%%%%

\subsection{Microplásticos}

Los microplásticos son fragmentos de plástico con un tamaño inferior a
5 mm.

Pueden originarse por:

\begin{itemize}

\item Fragmentación de objetos plásticos.

\item Desgaste de neumáticos.

\item Lavado de tejidos sintéticos.

\item Cosméticos y productos industriales.

\end{itemize}

Estos pequeños fragmentos pueden ser ingeridos por organismos acuáticos y entrar en
la cadena alimentaria.

%%%%%%%%%%%%%%%%%%%%%%%%%%%%%%%%%%%%%%%%%%%%%%%%%%%%%%%%%%%%

\subsection{Reciclaje de los polímeros}

El reciclaje constituye una de las estrategias más importantes para disminuir el
impacto ambiental de los residuos plásticos.

Existen tres procedimientos principales.

\subsubsection*{Reciclaje mecánico}

Consiste en triturar, fundir y volver a moldear el plástico.

Es el método más empleado para termoplásticos.

\vspace{0.3cm}

\subsubsection*{Reciclaje químico}

El polímero se transforma nuevamente en monómeros o en sustancias químicas de interés.

Permite obtener materiales de elevada calidad, aunque resulta más costoso.

\vspace{0.3cm}

\subsubsection*{Valorización energética}

Consiste en aprovechar la energía liberada durante la combustión controlada del
plástico.

Solo debe utilizarse cuando no sea posible reciclar el material.

%%%%%%%%%%%%%%%%%%%%%%%%%%%%%%%%%%%%%%%%%%%%%%%%%%%%%%%%%%%%

\subsection{Identificación de los plásticos}

Muchos envases incorporan un símbolo con un número que facilita su clasificación
durante el reciclaje.

\begin{center}

\renewcommand{\arraystretch}{1.3}

\begin{tabular}{ccl}

\toprule

\textbf{Código} &
\textbf{Polímero} &
\textbf{Ejemplo}

\\

\midrule

1 & PET & Botellas de bebidas \\

2 & PEAD & Garrafas y detergentes \\

3 & PVC & Tuberías y perfiles \\

4 & PEBD & Bolsas \\

5 & PP & Tapones y envases \\

6 & PS & Bandejas y vasos \\

7 & Otros & Policarbonato, PLA, etc. \\

\bottomrule

\end{tabular}

\end{center}

%%%%%%%%%%%%%%%%%%%%%%%%%%%%%%%%%%%%%%%%%%%%%%%%%%%%%%%%%%%%

\subsection{Bioplásticos y polímeros biodegradables}

Con el objetivo de reducir el impacto ambiental se están desarrollando nuevos
materiales obtenidos a partir de materias primas renovables.

Entre ellos destacan:

\begin{itemize}

\item PLA (ácido poliláctico).

\item PHA (polihidroxialcanoatos).

\item Plásticos obtenidos a partir de almidón.

\end{itemize}

Estos materiales pueden degradarse con mayor facilidad bajo determinadas condiciones,
aunque no todos son biodegradables en el medio natural.

%%%%%%%%%%%%%%%%%%%%%%%%%%%%%%%%%%%%%%%%%%%%%%%%%%%%%%%%%%%%

\subsection{Economía circular}

La gestión sostenible de los polímeros se basa en el modelo de economía circular,
cuyo objetivo es mantener los materiales en uso durante el mayor tiempo posible.

Las estrategias más importantes son las conocidas como las \textbf{5R}:

\begin{enumerate}

\item Rechazar.

\item Reducir.

\item Reutilizar.

\item Reparar.

\item Reciclar.

\end{enumerate}

Aplicar estas medidas permite disminuir el consumo de materias primas y reducir la
cantidad de residuos generados.

%%%%%%%%%%%%%%%%%%%%%%%%%%%%%%%%%%%%%%%%%%%%%%%%%%%%%%%%%%%%

\subsection{Resumen de los principales polímeros}

\begin{center}

\renewcommand{\arraystretch}{1.25}

\begin{tabular}{p{2.6cm}p{2.3cm}p{2cm}p{5cm}}

\toprule

\textbf{Polímero} &
\textbf{Tipo} &
\textbf{Obtención} &
\textbf{Aplicaciones}

\\

\midrule

PE &
Termoplástico &
Adición &
Bolsas, envases, tuberías

\\

PP &
Termoplástico &
Adición &
Tapones, recipientes

\\

PVC &
Termoplástico &
Adición &
Tuberías, cables

\\

PS &
Termoplástico &
Adición &
Embalajes

\\

PTFE &
Termoplástico &
Adición &
Recubrimientos antiadherentes

\\

PET &
Poliéster &
Condensación &
Botellas y fibras

\\

Nailon 6,6 &
Poliamida &
Condensación &
Tejidos y cuerdas

\\

Baquelita &
Termoestable &
Condensación &
Material eléctrico

\\

\bottomrule

\end{tabular}

\end{center}

%%%%%%%%%%%%%%%%%%%%%%%%%%%%%%%%%%%%%%%%%%%%%%%%%%%%%%%%%%%%

\begin{tcolorbox}[
title=Ideas clave,
colback=blue!5,
colframe=blue!70!black]

\begin{itemize}

\item Los polímeros son materiales indispensables en la sociedad actual.

\item Sus propiedades dependen de la estructura molecular y del tipo de polimerización.

\item El uso masivo de plásticos plantea importantes problemas ambientales.

\item El reciclaje, los bioplásticos y la economía circular constituyen estrategias
fundamentales para reducir su impacto.

\end{itemize}

\end{tcolorbox}

%%%%%%%%%%%%%%%%%%%%%%%%%%%%%%%%%%%%%%%%%%%%%%%%%%%%%%%%%%%%
%%%%%%%%%%%%%%%%%%%%%%%% SECCIÓN 8 %%%%%%%%%%%%%%%%%%%%%%%%%%
%%%%%%%%%%%%%%%%%%%%%%% RESUMEN FINAL %%%%%%%%%%%%%%%%%%%%%%%
%%%%%%%%%%%%%%%%%%%%%%%%%%%%%%%%%%%%%%%%%%%%%%%%%%%%%%%%%%%%

\section{Resumen}

Los polímeros son macromoléculas formadas por la unión repetitiva de unidades
más pequeñas llamadas \textbf{monómeros}. Gracias a su gran variedad de propiedades,
constituyen uno de los grupos de materiales más importantes de la industria moderna.

Dependiendo del mecanismo de formación, los polímeros se clasifican en dos grandes
grupos:

\begin{itemize}

\item \textbf{Polimerización por adición}: los monómeros, generalmente con dobles
enlaces, se unen sin desprender ninguna molécula.

\item \textbf{Polimerización por condensación}: la unión de los monómeros va
acompañada de la eliminación de una molécula pequeña, normalmente agua.

\end{itemize}

Las propiedades de un polímero dependen de:

\begin{itemize}

\item la naturaleza del monómero;

\item la longitud de las cadenas;

\item el grado de ramificación;

\item la reticulación;

\item las fuerzas intermoleculares.

\end{itemize}

Los polímeros presentan numerosas aplicaciones en la industria del envase, la
construcción, la medicina, el transporte y el sector textil. Sin embargo, su uso
masivo también plantea importantes problemas medioambientales relacionados con la
acumulación de residuos y la formación de microplásticos.

La reutilización, el reciclaje y el desarrollo de nuevos polímeros biodegradables
constituyen algunas de las principales estrategias para reducir su impacto sobre el
medio ambiente.

%%%%%%%%%%%%%%%%%%%%%%%%%%%%%%%%%%%%%%%%%%%%%%%%%%%%%%%%%%%%

\subsection{Mapa conceptual}

\begin{center}

\begin{tikzpicture}[
every node/.style={
draw,
rounded corners,
align=center,
font=\small
},
node distance=16mm
]

\node (A) {POLÍMEROS};

\node[below left=of A] (B)
{Clasificación};

\node[below=of A] (C)
{Polimerización};

\node[below right=of A] (D)
{Propiedades};

\node[below=of C] (E)
{Aplicaciones};

\node[right=of E] (F)
{Impacto\\ambiental};

\draw[-latex](A)--(B);
\draw[-latex](A)--(C);
\draw[-latex](A)--(D);
\draw[-latex](C)--(E);
\draw[-latex](E)--(F);

\end{tikzpicture}

\end{center}

%%%%%%%%%%%%%%%%%%%%%%%%%%%%%%%%%%%%%%%%%%%%%%%%%%%%%%%%%%%%

\subsection{Esquema general}

\begin{center}

\renewcommand{\arraystretch}{1.4}

\begin{tabular}{p{4cm}p{8cm}}

\toprule

\textbf{Aspecto}

&

\textbf{Idea clave}

\\

\midrule

Unidad básica

&

Monómero.

\\

Producto final

&

Polímero o macromolécula.

\\

Tipos de polimerización

&

Adición y condensación.

\\

Clasificación

&

Origen, composición, estructura y comportamiento térmico.

\\

Propiedades

&

Ligeros, resistentes y buenos aislantes.

\\

Problema ambiental

&

Residuos y microplásticos.

\\

Solución

&

Reducir, reutilizar, reciclar y desarrollar bioplásticos.

\\

\bottomrule

\end{tabular}

\end{center}

%%%%%%%%%%%%%%%%%%%%%%%%%%%%%%%%%%%%%%%%%%%%%%%%%%%%%%%%%%%%

\begin{tcolorbox}[
title=Ideas fundamentales para la EBAU,
colback=blue!5,
colframe=blue!70!black]

Para resolver correctamente los ejercicios sobre polímeros conviene recordar que:

\begin{enumerate}

\item En la polimerización por adición no se elimina ninguna molécula.

\item En la polimerización por condensación sí se elimina una molécula pequeña.

\item El nombre del polímero suele derivar del monómero del que procede.

\item La estructura molecular determina las propiedades del material.

\item El reciclaje y la economía circular son esenciales para disminuir el impacto
ambiental de los plásticos.

\end{enumerate}

\end{tcolorbox}

%%%%%%%%%%%%%%%%%%%%%%%%%%%%%%%%%%%%%%%%%%%%%%%%%%%%%%%%%%%%
%%%%%%%%%%%%%%%%%%%%%%% ANEXO %%%%%%%%%%%%%%%%%%%%%%%%%%%%%%%
%%%%%%%%%%%%%%%%%%%% TABLAS RESUMEN %%%%%%%%%%%%%%%%%%%%%%%%%
%%%%%%%%%%%%%%%%%%%%%%%%%%%%%%%%%%%%%%%%%%%%%%%%%%%%%%%%%%%%

\section*{Anexo. Tablas resumen}

\subsection*{1. Clasificación de los polímeros}

\begin{center}

\renewcommand{\arraystretch}{1.35}

\begin{tabular}{lll}

\toprule

\textbf{Criterio}

&

\textbf{Tipos}

&

\textbf{Ejemplos}

\\

\midrule

Origen

&

Naturales / Sintéticos

&

Celulosa, ADN / PE, PVC

\\

Composición

&

Homopolímeros / Copolímeros

&

PE / PET

\\

Estructura

&

Lineales / Ramificados / Reticulados

&

PEAD / PEBD / Baquelita

\\

Comportamiento térmico

&

Termoplásticos / Termoestables / Elastómeros

&

PET / Melamina / Caucho

\\

\bottomrule

\end{tabular}

\end{center}

%%%%%%%%%%%%%%%%%%%%%%%%%%%%%%%%%%%%%%%%%%%%%%%%%%%%%%%%%%%%

\subsection*{2. Tipos de polimerización}

\begin{center}

\renewcommand{\arraystretch}{1.35}

\begin{tabular}{p{4.3cm}p{4.2cm}p{4cm}}

\toprule

&

\textbf{Adición}

&

\textbf{Condensación}

\\

\midrule

Monómeros

&

Dobles enlaces

&

Grupos funcionales

\\

Pérdida de moléculas

&

No

&

Sí

\\

Productos secundarios

&

No

&

Sí

\\

Ejemplos

&

PE, PP, PVC, PS

&

PET, Nailon

\\

\bottomrule

\end{tabular}

\end{center}

%%%%%%%%%%%%%%%%%%%%%%%%%%%%%%%%%%%%%%%%%%%%%%%%%%%%%%%%%%%%

\subsection*{3. Principales polímeros}

\begin{center}

\scriptsize

\renewcommand{\arraystretch}{1.25}

\begin{tabular}{lllll}

\toprule

\textbf{Polímero}

&

\textbf{Monómero}

&

\textbf{Tipo}

&

\textbf{Propiedad}

&

\textbf{Aplicación}

\\

\midrule

PE

&

Eteno

&

Adición

&

Flexible

&

Bolsas

\\

PP

&

Propeno

&

Adición

&

Resistente

&

Envases

\\

PVC

&

Cloroeteno

&

Adición

&

Rígido

&

Tuberías

\\

PS

&

Estireno

&

Adición

&

Ligero

&

Embalajes

\\

PTFE

&

Tetrafluoroeteno

&

Adición

&

Antiadherente

&

Sartenes

\\

PET

&

Ácido tereftálico +
etanodiol

&

Condensación

&

Transparente

&

Botellas

\\

Nailon 6,6

&

Ácido adípico +
hexanodiamina

&

Condensación

&

Muy resistente

&

Fibras

\\

Baquelita

&

Fenol +
formaldehído

&

Condensación

&

Termoestable

&

Enchufes

\\

\bottomrule

\end{tabular}

\end{center}

%%%%%%%%%%%%%%%%%%%%%%%%%%%%%%%%%%%%%%%%%%%%%%%%%%%%%%%%%%%%

\subsection*{4. Las 5R}

\begin{center}

\Large

\textbf{
Rechazar
$\Longrightarrow$
Reducir
$\Longrightarrow$
Reutilizar
$\Longrightarrow$
Reparar
$\Longrightarrow$
Reciclar
}

\end{center}

%%%%%%%%%%%%%%%%%%%%%%%%%%%%%%%%%%%%%%%%%%%%%%%%%%%%%%%%%%%%

\begin{tcolorbox}[
title=Lo imprescindible para la EBAU,
colback=red!5,
colframe=red!70!black]

\begin{itemize}

\item Diferenciar claramente adición y condensación.

\item Conocer los monómeros de PE, PVC, PET y Nailon.

\item Saber representar la unidad repetitiva.

\item Relacionar estructura con propiedades.

\item Conocer las aplicaciones más importantes.

\item Explicar los problemas ambientales y las medidas para reducirlos.

\end{itemize}

\end{tcolorbox}